% --------------------------------------------------
% 7. CONCLUSION
% --------------------------------------------------
\section{Conclusion}

\begin{itemize}
    \item Restate the research question and summarise the main findings: previous DNF rate is the dominant predictor of abandonment; the Tour de France has a protective effect consistent with higher event value; mountain stages raise hazard consistent with marginal cost of effort; GC contenders show a counterintuitive positive hazard that reflects strategic withdrawal and physical exposure.
    \item Discuss limitations: no direct observation of financial variables (prize money, wages, contracts); the GC contender label is a coarse proxy; unobserved team-level strategy; potential selection bias in which riders start Grand Tours.
    \item Suggest extensions: incorporate prize-money or UCI-points data; add contract-year indicators; use competing-risks models to distinguish crash-related from fatigue-related DNFs; extend to one-week stage races for comparison.
    \item Closing sentence: tie back to the broader sports-economics literature and the value of using revealed-preference data to study economic decision-making under physical constraints.
\end{itemize}


% --------------------------------------------------
% REFERENCES
% --------------------------------------------------
\section*{References}
\addcontentsline{toc}{section}{References}

\begin{itemize}
    \item Kahn, L.\ M.\ (2000). The sports business as a labor market laboratory. \textit{Journal of Economic Perspectives}, 14(3), 75--94.
    \item Rosen, S.\ (1981). The economics of superstars. \textit{American Economic Review}, 71(5), 845--858.
    \item [Add course lecture references: F1--F7]
    \item [Add any additional references used]
\end{itemize}


% --------------------------------------------------
% APPENDIX
% --------------------------------------------------
\appendix
\section{Appendix}

\begin{itemize}
    \item Full regression output tables (if abridged in main text).
    \item Additional figures: Kaplan--Meier curves (Figure~1), DNF by stage type (Figure~3), cumulative DNF by stage number (Figure~4), Schoenfeld residual plots (Figure~6).
    \item Description of the web-scraping procedure and data-cleaning steps.
    \item Full counterfactual GC standings table (Table~6, if too long for main text).
\end{itemize}


% ============================================================
% REASONING FOR STRUCTURAL CHOICES (outcommented)
% ============================================================

\begin{comment}

WHY THIS STRUCTURE?

1. The exam document (KAN-CGMAV1006U) requires: Introduction, Chapters,
   Conclusion, References, and optionally Appendixes. This outline follows
   that structure exactly, with the "Chapters" expanded into four thematic
   sections (Theory, Data, Method, Results + Robustness).

2. The Introduction is designed to satisfy the exam requirement that it
   presents the research question, explains its importance, discusses the
   literature, material, and methods. All five elements are covered in
   the bullet points.

3. Theory/Literature is placed before Data and Method because the exam is
   for "The Economics of Sports" -- the graders want to see that you can
   connect your analysis to economic theory. By leading with tournament
   theory (Rosen 1981), labour market theory (Kahn 2000), and the
   team-as-firm framing from the lectures, you demonstrate command of the
   course material before ever showing a regression table.

4. The research question is deliberately framed around "proxy measures of
   economic incentives and continuation costs" rather than claiming to
   measure money directly. This is because the dataset contains no
   financial variables (no prize money, wages, or budgets). The framing
   is honest and defensible: you call your variables proxies, use careful
   language ("consistent with," "suggests"), and avoid overclaiming. This
   is the single most important thing the examiner will look for.

5. The Data section includes an explicit "economic interpretation" of
   each variable. This turns what could be a dry variable list into an
   economics argument: GC contender = tournament incentive, race dummies
   = event value, DNF rate = expected cost, mountain stage = marginal
   cost of effort. This mapping is the backbone of the paper's economics
   contribution.

6. The Empirical Strategy section is separate from Results to show
   methodological awareness. The Cox PH model is justified on its merits
   (semi-parametric, handles censoring, models duration). The
   time-varying extension is motivated by the economic intuition that
   stage difficulty varies within a race. The PH assumption test (Table 4)
   is included because it validates the model -- all p-values are above
   0.05, so the specification is sound.

7. Results are split into Basic Cox and Time-Varying Cox to show
   progression. The basic model gives the headline findings (DNF rate,
   GC contender, Tour de France effect). The time-varying model adds
   nuance (mountain stages matter, some rider-level effects are absorbed).
   An explicit "Economic Interpretation" subsection ties everything back
   to theory.

8. Robustness Checks serve two purposes:
   (a) The sub-sample analysis (GC vs non-GC) directly tests the
       tournament-incentive hypothesis by showing heterogeneous effects.
       The fact that DNF rate is much stronger for GC contenders (HR=2.28)
       than non-GC riders (HR=1.47) is a genuinely interesting economic
       finding.
   (b) The GC counterfactual analysis demonstrates the economic
       consequence of selective attrition -- it reshuffles prize money,
       UCI points, and performance signals. This moves the paper beyond
       "what predicts DNF" into "why does DNF matter economically."

9. The Conclusion is kept tight: summary, limitations, extensions. The
   limitations paragraph is critical because it shows self-awareness about
   what the data can and cannot say. Naming specific extensions (prize
   money data, contract-year indicators, competing risks) signals that
   you understand how the project could be improved.

10. Tables and figures are distributed across sections rather than dumped
    in an appendix. The exam is about analysis, and the graders want to
    see you discuss your tables. The appendix holds supplementary material
    and the full counterfactual table.

FIGURE AND TABLE MAPPING:
    Table 1 (summary stats)         -> Section 3.3
    Table 2 (basic Cox)             -> Section 5.1
    Table 3 (time-varying Cox)      -> Section 5.2
    Table 4 (Schoenfeld test)       -> Section 4.3
    Tables 5a/5b (subsamples)       -> Section 6.1
    Table 6 (counterfactual)        -> Section 6.2
    Figure 1 (KM curves)            -> Section 3.3 or Appendix
    Figure 2 (DNF by race/year)     -> Section 3.3
    Figure 3 (DNF by stage type)    -> Appendix
    Figure 4 (cumulative DNF)       -> Appendix
    Figure 5 (hazard ratio forest)  -> Section 5.1
    Figure 6 (Schoenfeld residuals) -> Section 4.3 or Appendix
    Figure 7 (GC vs non-GC forest)  -> Section 6.1
    Figure 8 (counterfactual GC)    -> Section 6.2

\end{comment}